Every package in this suite was built by an AI coding agent, every package passed the checks its builder wrote, and an LLM council approved the pull requests that shipped them. An independent review that probed the built packages still found defects in all \suitePackages{} of them. That record is the finding of this paper: agents can build benchmarks, and approval is not review.

The artifact underneath the finding is small. \texttt{benchmark-core} is a shared TypeScript engine. \suitePackages{} packages sit on it, each with its own tasks and deterministic grader, published on npm and under \path{github.com/pooriaarab/benchmark-*}. One \texttt{defineBenchmark} call produces four ways to run the same task set: an SDK import, a command-line interface, an MCP server, and a fetch-style HTTP handler. A mock provider runs the whole suite offline with no key, so a second machine can check every grader without a model. Figure~\ref{fig:system} shows the shape. The directory that lists the packages is \url{https://private-repo-337e.com}.

The review record is public. All \reviewPackages{} packages needed at least one fix round, meaning one cycle of independent findings and their fixes (median \reviewRoundsMedian{}, maximum \reviewRoundsMax{}; Figure~\ref{fig:rounds}). The first \reviewWaveOne{} shipped a 0.2.0 whose changelog states that 0.1.0 scores are not comparable. At least \defectBlindStrategyCredit{} packages record a strategy that ignored the task and still earned credit. Section~\ref{sec:review} names \defectClasses{} defect classes, counts the packages each one touched, and cites the issue, pull request, or changelog line behind every cell (Table~\ref{tab:defects}, Appendix~\ref{app:locators}).

The approvals are on the record too, and they are the counter-evidence this paper has to carry. The council workflow that reviewed the builder pull requests posted \councilReviews{} reviews over \councilPrsReviewed{} of the \buildersPrs{} pull requests in the \suitePackages{} repositories. It approved pelican's grader pull request; an issue in the same repository later records that one memorised SVG scored 1.0 on every task. The defect findings came from separate review sessions that attacked the packages under a probing protocol (Appendix~\ref{app:protocol}). The claim the data supports is about the depth of review, not about who reviewed what: green checks and approvals did not find what adversarial probing found. The builders were never asked to review their own work, the probing reviewers were agents as well, there is no human-built control, the suite is one pipeline of \suitePackages{} packages, and no real model has been run, so this paper claims nothing about any model's ability.

Our contributions are:
\begin{enumerate}
\item the defect finding and a sourced taxonomy of \defectClasses{} classes across \suitePackages{} agent-built packages (Section~\ref{sec:review}, Table~\ref{tab:defects}, Figure~\ref{fig:defects});
\item the per-package fix-round record (Figure~\ref{fig:rounds}) and the re-review protocol the independent sessions followed, short enough to reuse (Appendix~\ref{app:protocol});
\item the artifact: one engine, \suitePackages{} packages, four surfaces each, with provenance and a licence per task source (Sections~\ref{sec:engine} and \ref{sec:suite});
\item validity documentation per package: construct, scope, and task generation in one table (Section~\ref{sec:validity});
\item the checks we could run without a model: oracle-versus-trivial separation with seeded bootstrap intervals, and a plain account of what those checks do not show (Section~\ref{sec:validation}).
\end{enumerate}

The pelican package needs a credit up front. The prompt ``Generate an SVG of a pelican riding a bicycle'' is Simon Willison's, used since October 2024 as a quick look at whether a model can draw~\cite{willison2024pelican}. Our package turns that prompt into a structural grader and a seeded grid of variants.
